\chapter{Podsumowanie}

Niniejsza praca przedstawiła kompleksowe rozwiązanie problemu autoremediacji systemów mikrousługowych opartego na dynamicznie ewoluujących kontekstach agentowych. Przeprowadzone eksperymenty na dwóch różnych modelach językowych - GPT-5-nano oraz Minimax-2.5 - ujawniły fundamentalne zależności między architekturą modelu a skutecznością potoku Agentic Context Engineering (ACE), które znacząco wpływają na projektowanie przyszłych systemów wieloagentowych.

\section{Główne osiągnięcia i wnioski}

\subsection{Potwierdzenie głównej hipotezy z istotnymi modyfikacjami}

Eksperyment w pełni potwierdził główną hipotezę pracy: \textbf{dynamiczna ewolucja kontekstów agentowych istotnie podnosi skuteczność autoremediacji}. Jednak kluczowym odkryciem jest silna zależność tej skuteczności od architektury wykorzystywanego modelu językowego (tabela~\ref{tab:model-comparison}):

\begin{itemize}
    \item \textbf{GPT-5-nano}: Dramatyczna poprawa o +18,4 punktu procentowego (+120\% względnie) we wskaźniku skuteczności remediacji przy niskiej skuteczności bazowej (15,3\%) - szczegóły w tabeli~\ref{tab:overall-success-gpt5}.
    \item \textbf{Minimax-2.5}: Umiarkowana poprawa o +5,1 punktu procentowego (+13\% względnie) przy wyższej skuteczności bazowej (38,8\%) - szczegóły w tabeli~\ref{tab:overall-success-minimax}.
\end{itemize}

Te wyniki wskazują, że modele o wyższych zdolnościach rozumowania (GPT-5-nano z \texttt{reasoning\_effort=medium}) lepiej absorpcją i wykorzystują strukturizowaną wiedzę proceduralną zawartą w ewolucyjnych Playbookach.

\subsection{Odkrycie model-specyficznych strategii organizacji pracy}

Analiza korelacji między strukturą drzewa zadań a sukcesem ujawniła fundamentalne różnice w strategiach organizacyjnych:

\paragraph{GPT-5-nano: Strategia delegacji cyrkulacyjnej}
Model wykazuje pozytywne korelacje zarówno dla głębokości (r=0.316***, p<0.001) jak i szerokości (r=0.351***, p<0.001) drzewa zadań (tabela~\ref{tab:correlation-analysis}, rysunek~\ref{fig:fig6-tree-structure-gpt5}). Analiza n-gramowa z korektą liftu ujawnia, że sukces koreluje z przejściami między różnymi agentami (\texttt{monitoring→coder}, lift~1,55x; \texttt{devops→monitoring}, lift~1,46x), a ponowna delegacja do tego samego agenta jest predyktorem porażki - szczegóły w podrozdziale~\ref{subsec:wzorce-gpt}. Każda kolejna delegacja to statystycznie kolejna zmiana agenta, co wyjaśnia dodatnią korelację szerokości.

\paragraph{Minimax-2.5: Uczenie się zamknięcia cyklu remediacji}
Model potwierdza hipotezę o pozytywnej korelacji głębokości (r=0.143*, p=0.045), ale wykazuje ujemną korelację szerokości (r=-0.168*, p=0.019) (tabela~\ref{tab:correlation-analysis}, rysunek~\ref{fig:fig6-tree-structure-minimax}). Jedynym istotnym wzorcem delegacyjnym jest \texttt{coder→monitoring} (lift~1,50x, $p<0{,}001$). Wzorce sukcesu leżą w sekwencjach narzędziowych prowadzących do wdrożenia (\texttt{update\_todo~$\to$~deploy\_app}, lift~1,68x), lecz sekwencja ta stanowi mechaniczny artefakt każdego udanego deploymenta, nie wyróżniającą się strategię organizacyjną. Faktyczna zmiana wprowadzona przez ACE to eliminacja pętli diagnostycznych monitoring$\leftrightarrow$devops (lift~$<0{,}2$) i nauka zamykania cyklu remediacji przez wdrożenie - szczegóły w podrozdziale~\ref{subsec:wzorce-minimax}. Asymetria między modelami jest głębsza niż różnica preferowanych strategii: GPT-5-nano uczy się \emph{jak koordynować} agentów, Minimax-2.5 uczy się \emph{kiedy wdrożyć}.


\section{Odpowiedzi na pytania badawcze}

Na podstawie przeprowadzonych eksperymentów można udzielić precyzyjnych odpowiedzi na pytania badawcze postawione w Rozdziale~I:

\subsection{Pytanie 1: Wpływ ewolucji Playbooka na skuteczność per klasa incydentu}

\textbf{Odpowiedź}: \textbf{Wpływ jest zróżnicowany między modelami. GPT-5-nano poprawia skuteczność we wszystkich czterech klasach; Minimax-2.5 poprawia klasy~2, 3 i~4, lecz wykazuje regresję w klasie~1 (usterki hybrydowe).}

Szczegółowe wyniki per klasa przedstawiono na rysunkach~\ref{fig:learning-comparison-gpt5} i~\ref{fig:learning-comparison-minimax} oraz w tabelach~\ref{tab:by-incident-type-gpt5} i~\ref{tab:by-incident-type-minimax}. GPT-5-nano poprawia skuteczność we wszystkich czterech klasach, przy czym największy przyrost odnotowuje w Klasie~1 ($+0{,}42$ pkt, $+76{,}9\%$) i Klasie~3 ($+0{,}29$ pkt, $+41{,}7\%$). Minimax-2.5 poprawia klasy~2, 3 i~4, lecz wykazuje regresję w Klasie~1 ($-0{,}08$ pkt, $-13{,}3\%$).

\subsection{Pytanie 2: Profil zmian charakterystyk operacyjnych}

\textbf{Odpowiedź}: \textbf{Profile operacyjne obu modeli różnią się jakościowo - GPT-5-nano płaci wysoką cenę w TTR, lecz uzyskuje dramatyczny wzrost skuteczności; Minimax-2.5 poprawia większość metryk przy umiarkowanym zysku efektywności.}

\paragraph{GPT-5-nano} (rysunek~\ref{fig:fig2-operational-gpt5})
\begin{itemize}
    \item \textbf{TTR}: wzrost mediany z 3,2 do 14,2~min ($+439{,}7\%$) - efekt selekcji: bez uczenia model szybko kończy epizod porażką, z uczeniem podejmuje głębszą eksplorację i częściej dochodzi do naprawy.
    \item \textbf{Liczba iteracji}: wzrost z 5,4 do 8,5 ($+56{,}2\%$).
    \item \textbf{Wywołania narzędzi}: wzrost z 7,7 do 17,2 ($+123{,}8\%$).
    \item \textbf{Łączna liczba błędów narzędziowych}: wzrost z 362 do 1\,843 ($+409{,}1\%$).
    \item \textbf{Struktura drzewa zadań}: głębokość rośnie z 0,86 do 1,01 ($+17{,}9\%$), liczba potomków z 3,2 do 4,8 ($+53{,}7\%$). Obie metryki korelują pozytywnie ze skutecznością remediacji: głębokość r$=0{,}316^{***}$ ($p{<}0{,}001$), szerokość r$=0{,}351^{***}$ ($p{<}0{,}001$) - każda dodatkowa delegacja to statystycznie kolejna zmiana agenta, co sprzyja sukcesowi.
    \item \textbf{Kompromis}: przy wzroście skuteczności o $+120\%$ wysoki koszt operacyjny jest uzasadniony - system uczy się eksplorować zamiast szybko kapitulować.
\end{itemize}

\paragraph{Minimax-2.5} (rysunek~\ref{fig:fig2-operational-minimax})
\begin{itemize}
    \item \textbf{TTR}: spadek mediany z 11,1 do 8,9~min ($-20{,}4\%$) - Playbook przyspiesza diagnozę i eliminuje pętle diagnostyczne.
    \item \textbf{Liczba iteracji}: wzrost z 15,4 do 24,9 ($+61{,}0\%$).
    \item \textbf{Wywołania narzędzi}: spadek z 22,1 do 18,8 ($-15{,}0\%$) - precyzyjniejsze kroki zastępują redundantne zapytania.
    \item \textbf{Łączna liczba błędów narzędziowych}: spadek z 2\,625 do 1\,788 ($-31{,}9\%$) - wyraźna poprawa jakości wywołań.
    \item \textbf{Struktura drzewa zadań}: głębokość rośnie z 1,10 do 1,39 ($+25{,}9\%$), liczba potomków z 4,4 do 5,4 ($+23{,}2\%$). Korelacja głębokości ze skutecznością jest słabo pozytywna (r$=0{,}143^{*}$, $p{=}0{,}045$), natomiast korelacja szerokości jest paradoksalnie ujemna (r$=-0{,}168^{*}$, $p{=}0{,}019$) - dodatkowe delegacje zwiększają ryzyko pętli diagnostycznych, co obniża skuteczność.
    \item \textbf{Kompromis}: przy wzroście skuteczności o $+13\%$ redukcja TTR i błędów narzędziowych wskazuje na korzystną zmianę profilu operacyjnego, lecz marginalny zysk efektywności może nie uzasadniać kosztów utrzymania potoku ACE dla tego modelu.
\end{itemize}

\subsection{Pytanie 3: Emergentne wzorce organizacji pracy}

\textbf{Odpowiedź}: \textbf{System wykształca model-specyficzne wzorce hierarchicznej delegacji, których charakter bezpośrednio wyjaśnia obserwowane korelacje struktury drzewa zadań z sukcesem.}

Najistotniejszą zmianą strukturalną jest pogłębienie drzewa na poziomie~2 (rysunki~\ref{fig:fig6-tree-structure-gpt5} i~\ref{fig:fig6-tree-structure-minimax}, rysunek~\ref{fig:przyklad-rozbudowanego-drzewa}): średnia liczba dzieci na tym poziomie wzrasta z 0,2 do 0,8, a odsetek epizodów z zadaniami na poziomie~2 rośnie z 13\% do 30\%. Agenci pośredniego szczebla sami stają się delegatorami - bez jawnego zaprogramowania tej strategii.

Analiza korelacji (tabela~\ref{tab:correlation-analysis}) i n-gramów (podrozdziały~\ref{subsec:wzorce-gpt}--\ref{subsec:interpretacja}) ujawnia dwa jakościowo odmienne wzorce:

\begin{itemize}
    \item \textbf{GPT-5-nano - delegacja cyrkulacyjna}: pozytywne korelacje zarówno głębokości (r=0,316***, p$<$0,001) jak i szerokości (r=0,351***, p$<$0,001). Sukces koreluje z przejściami między różnymi agentami (\texttt{monitoring$\to$coder}, lift~1,55x; \texttt{devops$\to$monitoring}, lift~1,46x), a ponowna delegacja do tego samego agenta jest predyktorem porażki (lift~0,00x). Każda kolejna delegacja to statystycznie kolejna zmiana roli - stąd dodatnia korelacja szerokości. Analiza per klasie ujawnia jednak, że \texttt{devops$\to$monitoring} (lift ogólny 1,46, GPT-5-nano) koncentruje się w Klasie~3 (lift 2,43, $p=0{,}003^{**}$) --- wzorzec jest klasa-specyficzny, nie globalny.
    \item \textbf{Minimax-2.5 - zamknięcie cyklu remediacji}: pozytywna korelacja głębokości (r=0,143*, p=0,045), lecz ujemna korelacja szerokości (r=$-0{,}168$*, p=0,019). ACE eliminuje pętle diagnostyczne monitoring$\leftrightarrow$devops (lift~$<0{,}2$) i uczy model doprowadzania naprawy do etapu wdrożenia zamiast przerywania jej w fazie diagnozy. Każda dodatkowa delegacja wiąże się z ryzykiem wpadnięcia w pętlę - stąd ujemna korelacja szerokości.
\end{itemize}

Dynamikę emergencji tych wzorców (porównanie L vs NL oraz ewolucja L1$\to$L2) opisano w podrozdziale~\ref{subsec:emergencja-dynamika}. Asymetria jest głębsza niż różnica preferowanych strategii: GPT-5-nano uczy się \emph{jak koordynować} agentów, Minimax-2.5 uczy się \emph{kiedy wdrożyć}.

\subsection{Pytanie 4: Bilans progresji i regresji oraz podatność klas incydentów na negatywny wpływ ACE}

\textbf{Odpowiedź}: \textbf{W obu modelach progresje przewyższają regresje, jednak Klasa~1 (usterki hybrydowe) jest szczególnie podatna na negatywny wpływ ewolucji Playbooka - wyłącznie w modelu Minimax-2.5.}

W wariancie z aktywnym potokiem ACE (tabele~\ref{tab:progressions-gpt5} i~\ref{tab:progressions-minimax}):
\begin{itemize}
    \item \textbf{GPT-5-nano}: 36 progresji (36,7\%) wobec 23 regresji (23,5\%) - stosunek $\approx 1{,}57$:1.
    \item \textbf{Minimax-2.5}: 39 progresji (39,8\%) wobec 23 regresji (23,5\%) - stosunek $\approx 1{,}70$:1.
\end{itemize}

Klasą podatną na regresję jest Klasa~1 (usterki hybrydowe) - ale wyłącznie dla Minimax-2.5 ($-0{,}08$ pkt, $-13{,}3\%$, tabela~\ref{tab:by-incident-type-minimax}). Źródłem regresji są heurystyki \texttt{kubectl\_exec} zapisane w Playbooku \textit{DevOps Agenta}, które okazują się kontrproduktywne. Uzupełnieniem jest wzorzec narzędziowy \texttt{deploy\_app $\to$ assign\_task} (Minimax-2.5, Klasa~1, lift 4,80, $p=0{,}005^{**}$, $n=5$): wysokie lift charakteryzuje rzadkie epizody sukcesów, ale zawodność tej strategii w niestabilnych podach odpowiada za regresję globalną klasy. GPT-5-nano nie wykazuje regresji w żadnej klasie, co wskazuje na model-specyficzną podatność - zdolność do selektywnego stosowania heurystyk decyduje o tym, czy nowe reguły działają w odpowiednich kontekstach. Kierunkiem dalszych prac jest warunkowanie reguł Playbooka wykrytą klasą incydentu (zob. podrozdział~\ref{sec:ewolucja-playbooka}).

\section{Znaczenie statystyczne wyników}

\subsection{Potwierdzone zależności statystyczne}

Eksperyment wykazał następujące statystycznie istotne zależności (tabela~\ref{tab:correlation-analysis} i \ref{tab:model-comparison}):

\begin{itemize}
    \item \textbf{GPT-5-nano z uczeniem}: Wszystkie korelacje struktury drzewa z sukcesem są statystycznie istotne (p<0.001)
    \item \textbf{Minimax-2.5 z uczeniem}: Słabsze, ale istotne korelacje głębokości (p=0.045) i ujemne korelacje szerokości (p=0.019)
    \item \textbf{Oba modele bez uczenia}: Brak statystycznie istotnych korelacji (p>0.05)
    \item \textbf{Poprawa wskaźnika skuteczności remediacji ($\ge$2)}: Istotna dla GPT-5-nano (p=0.005) i Minimax-2.5 (p=0.562)
\end{itemize}

\subsection{Obserwowane różnice między modelami}

Wyniki eksperymentu ujawniły znaczące różnice w zachowaniu modeli (tabela~\ref{tab:model-comparison} i~\ref{tab:correlation-analysis}):
\begin{itemize}
    \item \textbf{Skuteczność bazowa}: Minimax-2.5 (38.8%) > GPT-5-nano (15.3%) bez uczenia
    \item \textbf{Potencjał uczenia}: GPT-5-nano (+18.4pp) > Minimax-2.5 (+5.1pp)
    \item \textbf{Korelacja szerokości}: GPT-5-nano pozytywna vs Minimax-2.5 negatywna
\end{itemize}

\section{Podsumowanie końcowe}

Eksperymenty przeprowadzone na dwóch modelach językowych (GPT-5-nano i Minimax-2.5) wykazały, że potok Agentic Context Engineering poprawia skuteczność autoremediacji w obu przypadkach, jednak w różnym stopniu. 

Kluczowym obserwacją jest różnica w korelacjach struktury drzewa zadań z sukcesem: dla GPT-5-nano wszystkie korelacje są pozytywne i silnie istotne statystycznie, podczas gdy dla Minimax-2.5 korelacja szerokości jest negatywna. To wskazuje na różne optymalne strategie organizacji pracy dla różnych architektury modeli językowych.

Przedstawiony potok ACE stanowi funkcjonalne rozwiązanie dla dynamicznej ewolucji kontekstów agentowych, z zastrzeżeniem konieczności uwzględnienia specyfiki wykorzystywanego modelu językowego.